% GENERATED FILE. Edit canonical lessons, then run npm run generate:books.
% canonical-source-hash: fnv1a64:cb56c4caf074b1ec
% canonical-lessons: JA-C48-niwatori, JA-C48-chiisai, JA-C48-ii, JA-C48-nagai, JA-C48-mijikai

\chapter{Chicken, Small, Good, Long, Short}
\label{ch:ja-chicken-small-good-long-short}
\hlchaptermodality{48}

\begin{chapteropening}
\textbf{By the end of this chapter:} \emph{I can say a chicken, small, good, long and short.}

Complete the last lesson of chapter 48: I can say a chicken, small, good, long and short.
\end{chapteropening}

\section[\texorpdfstring{\ja{にわとり} (niwatori)}{にわとり (niwatori)}]{\ja{にわとり} (niwatori) — a chicken}

\label{lesson:JA-C48-niwatori}



\begin{quote}
Before the new one: say the Japanese for a turtle, then the Japanese for a wolf.
\end{quote}

\subsection*{You'll want to know: \ja{にわとり}}
\textbf{\ja{にわとり}} — \emph{niwatori} — \textquotedblleft{}a chicken\textquotedblright{}.

\begin{cousinweb}
\textbf{\ja{にわとり}} (\emph{niwatori}) is a native compound, inherited rather than borrowed: \emph{niwa}, \textquotedblleft{}yard, garden\textquotedblright{}, and \emph{tori}, \textquotedblleft{}bird\textquotedblright{}, the bird of the yard. JMdict writes it with one character for \textquotedblleft{}chicken\textquotedblright{}, read as a whole.
\end{cousinweb}

\begin{grammarlens}[title={What you've built}]
One more word for this chapter.
\end{grammarlens}

\subsection*{Your turn}
\begin{itemize}
\raggedright
  \item \emph{Say it:} \emph{niwatori}
  \item \emph{Say it:} \emph{niwatori}, once more
  \item \emph{Say it:} say \emph{niwatori}
  \item \emph{Recall:} say the Japanese for a turtle, then the Japanese for a wolf, then say \emph{niwatori} again
\end{itemize}

\begin{tcolorbox}[breakable,colback=teal!4,colframe=teal!35!black,title={Before you move on}]
What does \ja{にわとり} mean? (\textquotedblleft{}A chicken\textquotedblright{}.) Say it once more.
\end{tcolorbox}
\section[\texorpdfstring{\ja{ちいさい} (chiisai)}{ちいさい (chiisai)}]{\ja{ちいさい} (chiisai) — small}

\label{lesson:JA-C48-chiisai}



\begin{quote}
Before the new one: say the Japanese for a wolf, then the Japanese for a chicken.
\end{quote}

\subsection*{You'll want to know: \ja{ちいさい}}
\textbf{\ja{ちいさい}} — \emph{chiisai} — \textquotedblleft{}small\textquotedblright{}.

\begin{cousinweb}
\textbf{\ja{ちいさい}} (\emph{chiisai}) is a native Japanese word, inherited rather than borrowed from Chinese, and its older origin is not settled. Its kanji was attached to it later as a kun'yomi, a native reading given to a Chinese character (JMdict, KANJIDIC2).
\end{cousinweb}

\begin{grammarlens}[title={What you've built}]
One more word for this chapter.
\end{grammarlens}

\subsection*{Your turn}
\begin{itemize}
\raggedright
  \item \emph{Say it:} \emph{chiisai}
  \item \emph{Say it:} \emph{chiisai}, once more
  \item \emph{Say it:} say \emph{chiisai}
  \item \emph{Recall:} say the Japanese for a wolf, then the Japanese for a chicken, then say \emph{chiisai} again
\end{itemize}

\begin{tcolorbox}[breakable,colback=teal!4,colframe=teal!35!black,title={Before you move on}]
What does \ja{ちいさい} mean? (\textquotedblleft{}Small\textquotedblright{}.) Say it once more.
\end{tcolorbox}
\section[\texorpdfstring{\ja{いい} (ii)}{いい (ii)}]{\ja{いい} (ii) — good}

\label{lesson:JA-C48-ii}



\begin{quote}
Before the new one: say the Japanese for a chicken, then the Japanese for small.
\end{quote}

\subsection*{You'll want to know: \ja{いい}}
\textbf{\ja{いい}} — \emph{ii} — \textquotedblleft{}good\textquotedblright{}.

\begin{cousinweb}
\textbf{\ja{いい}} (\emph{ii}) is a native word, a later spoken form of \emph{yoi}, \textquotedblleft{}good\textquotedblright{}; JMdict notes that \emph{ii} stands for \emph{yoi} in all its senses. It is inherited rather than borrowed.
\end{cousinweb}

\begin{grammarlens}[title={What you've built}]
One more word for this chapter.
\end{grammarlens}

\subsection*{Your turn}
\begin{itemize}
\raggedright
  \item \emph{Say it:} \emph{ii}
  \item \emph{Say it:} \emph{ii}, once more
  \item \emph{Say it:} say \emph{ii}
  \item \emph{Recall:} say the Japanese for a chicken, then the Japanese for small, then say \emph{ii} again
\end{itemize}

\begin{tcolorbox}[breakable,colback=teal!4,colframe=teal!35!black,title={Before you move on}]
What does \ja{いい} mean? (\textquotedblleft{}Good\textquotedblright{}.) Say it once more.
\end{tcolorbox}
\section[\texorpdfstring{\ja{ながい} (nagai)}{ながい (nagai)}]{\ja{ながい} (nagai) — long}

\label{lesson:JA-C48-nagai}



\begin{quote}
Before the new one: say the Japanese for small, then the Japanese for good.
\end{quote}

\subsection*{You'll want to know: \ja{ながい}}
\textbf{\ja{ながい}} — \emph{nagai} — \textquotedblleft{}long\textquotedblright{}.

\begin{cousinweb}
\textbf{\ja{ながい}} (\emph{nagai}) is a native Japanese word, inherited rather than borrowed from Chinese, and its older origin is not settled. Its kanji was attached to it later as a kun'yomi, a native reading given to a Chinese character (JMdict, KANJIDIC2).
\end{cousinweb}

\begin{grammarlens}[title={What you've built}]
One more word for this chapter.
\end{grammarlens}

\subsection*{Your turn}
\begin{itemize}
\raggedright
  \item \emph{Say it:} \emph{nagai}
  \item \emph{Say it:} \emph{nagai}, once more
  \item \emph{Say it:} say \emph{nagai}
  \item \emph{Recall:} say the Japanese for small, then the Japanese for good, then say \emph{nagai} again
\end{itemize}

\begin{tcolorbox}[breakable,colback=teal!4,colframe=teal!35!black,title={Before you move on}]
What does \ja{ながい} mean? (\textquotedblleft{}Long\textquotedblright{}.) Say it once more.
\end{tcolorbox}
\section[\texorpdfstring{\ja{みじかい} (mijikai)}{みじかい (mijikai)}]{\ja{みじかい} (mijikai) — short}

\label{lesson:JA-C48-mijikai}



\begin{quote}
Before the new one: say the Japanese for good, then the Japanese for long.
\end{quote}

\subsection*{You'll want to know: \ja{みじかい}}
\textbf{\ja{みじかい}} — \emph{mijikai} — \textquotedblleft{}short\textquotedblright{}.

\begin{cousinweb}
\textbf{\ja{みじかい}} (\emph{mijikai}) is a native Japanese word, inherited rather than borrowed from Chinese, and its older origin is not settled. Its kanji was attached to it later as a kun'yomi, a native reading given to a Chinese character (JMdict, KANJIDIC2).
\end{cousinweb}

\begin{grammarlens}[title={What you've built}]
One more word for this chapter.
\end{grammarlens}

\subsection*{Your turn}
\begin{itemize}
\raggedright
  \item \emph{Say it:} \emph{mijikai}
  \item \emph{Say it:} \emph{mijikai}, once more
  \item \emph{Say it:} say \emph{mijikai}
  \item \emph{Recall:} say the Japanese for good, then the Japanese for long, then say \emph{mijikai} again
\end{itemize}

\begin{tcolorbox}[breakable,colback=teal!4,colframe=teal!35!black,title={Before you move on}]
What does \ja{みじかい} mean? (\textquotedblleft{}Short\textquotedblright{}.) Say it once more.
\end{tcolorbox}
